% Source PDF pages 30--32; manual pages 2-1--2-3.
As mentioned in the previous section, TAOS uses the same methods as its predecessor
TSAP,\taosref{2} so the description of the mathematical formulation given in the
TSAP manual applies to TAOS as well. However, that description is incomplete. This
chapter includes the TSAP material together with additional information on coordinate
systems, equations of motion, atmosphere and gravity models, guidance methods,
range calculations, searches, and optimization.

\taossubhead{Notation}

Although a list of symbols is included in the nomenclature, some of the notation
requires additional explanation. TAOS uses many coordinate systems, so the notation
is designed to state coordinate-system relationships precisely.

Vectors of arbitrary magnitude, used for position, velocity, and acceleration, are
denoted with an arrow. Vector magnitude is denoted with double bars, unit vectors
with a caret, and matrices with square brackets:
\[
\begin{aligned}
  \vec{x}       &={} \text{vector of arbitrary magnitude},\\
  \lVert\vec{x}\rVert &={} \text{magnitude of vector},\\
  \hat{x}       &={} \text{unit vector},\\
  [A]           &={} \text{matrix}.
\end{aligned}
\]

A right subscript on a vector indicates a coordinate-system association. For unit
vectors, it identifies the coordinate system represented by the unit vector. For other
vectors, it identifies the coordinate-system dependence produced by differentiation
with respect to time. For example, $\vec{V}_{\oplus}$ is the velocity vector relative
to an earth-fixed coordinate system.

Left and right subscripts indicate that a value is relative to two reference frames.
Thus, ${}_{I}\vec{\omega}_{\oplus}$ is the rotational velocity between the inertial
and earth-fixed coordinate systems. This notation is also used for transformation
matrices and other quantities that convert from one system to another.

Derivatives with respect to time are denoted with one or two dots above the symbol:
\[
\begin{aligned}
  \dot{h}  &={} \text{first derivative of $h$ with respect to time},\\
  \ddot{h} &={} \text{second derivative of $h$ with respect to time}.
\end{aligned}
\]

The dot product between two vectors can be interpreted as the projection of one
vector onto another. Projecting a vector onto the unit vectors of a coordinate system
gives the components of the vector in that system. For example,
\[
\begin{aligned}
  \vec{V}_{\oplus}\mathbin{\cdot}\hat{x}_{gd}
    &={} \text{component of earth-fixed velocity in the geodetic $x$ direction},\\
  \vec{V}_{\oplus}\mathbin{\cdot}\hat{x}_{b}
    &={} \text{component of earth-fixed velocity in the body $x$ direction}.
\end{aligned}
\]

\taossubhead{Numerical Errors}

The equations of motion are purposely written in a nonsingular earth-fixed form to
avoid numerical problems during integration, such as division by zero. The only two
conditions that cause the equations of motion themselves to fail are a vehicle at the
center of the earth and a vehicle with zero mass. If either condition occurs, the
trajectory is ended and an error message is issued.

Although the equations of motion are nonsingular, many output-variable equations
contain denominators that can become zero in special situations. If this occurs,
TAOS assumes a reasonable value, such as zero, and continues execution.

\taossubhead{Program Modules}

The TAOS source code is divided into modules containing related functions or
subroutines, as shown in \cref{tab:taos-program-modules}.

\begin{table}[htbp]
  \centering
  \caption{TAOS program modules.}
  \label{tab:taos-program-modules}
  \begin{tabularx}{0.68\linewidth}{@{}>{\ttfamily}lX@{}}
    \toprule
    \normalfont Module & Description \\
    \midrule
    atmos  & Model atmospheres \\
    coords & Coordinate systems \\
    derivs & Equations of motion and guidance \\
    errors & Error messages \\
    in\_prb & Problem input \\
    in\_seg & Trajectory-segment input \\
    in\_trj & Trajectory input \\
    main   & Main program \\
    numeric & Numerical analysis \\
    output & Output files \\
    path   & Trajectory calculation and control \\
    search & Search and optimization control \\
    tables & Table input and interpolation \\
    util   & Utility functions \\
    vf02   & Optimization \\
    \bottomrule
  \end{tabularx}
\end{table}

Many of these modules, such as \texttt{errors}, \texttt{in\_prb}, \texttt{output},
and \texttt{tables}, perform input or output functions and are documented in
\cref{chap:table-files,chap:problem-files}. The modules containing trajectory
calculations---\texttt{atmos}, \texttt{coords}, \texttt{derivs}, \texttt{numeric},
\texttt{path}, and \texttt{search}---are documented in this chapter.

\taossubhead{Units}

Internally, TAOS uses English units for everything except atmosphere calculations,
which use metric units. The input and output units can be changed from the default
English system to metric units with the \problemblock{*Units/Fmt} data block
(\cref{sec:block-units-fmt}).

The following conversion factors are used in TAOS:
\begin{center}
\begin{tabular}{@{}r l@{}}
  \num{0.017453293} & radians per degree \\
  \num{57.29577951} & degrees per radian \\
  \num{2.204622476} & pounds mass per kilogram \\
  \num{0.45359240} & kilograms per pound mass \\
  \num{32.1740485} & pounds mass per slug \\
  \num{0.03108095} & slugs per pound mass \\
  \num{3.280839895} & feet per meter \\
  \num{0.3048} & meters per foot \\
  \num{6076.1} & feet per nautical mile \\
  \num{0.0001645792} & nautical miles per foot \\
  \num{5280.0} & feet per statute mile \\
  \num{0.0001893939} & statute miles per foot \\
  \num{0.224808924} & pounds force per newton \\
  \num{4.448222} & newtons per pound force \\
  \num{0.020885434} & pounds per square foot per pascal \\
  \num{47.88026} & pascals per pound per square foot \\
  \num{9.806650} & acceleration of gravity in meters per second squared \\
\end{tabular}
\end{center}

\noindent\textit{Editorial note.} The descriptions of the first two angular
conversion factors are reversed in the printed source; the mathematically consistent
labels are used here.
